% -------------------------------------------------------------------------
%  fig-add.tex  --  the five steps of add()
% -------------------------------------------------------------------------
\newcommand{\FigAdd}{%
\begin{figure}[tb]
\centering
\begin{tikzpicture}[x=1mm,y=1mm]

  \node[lbl, anchor=west] at (-16,5.5) {\itshape chain stored at $\dk$ before the insertion};
  \foreach \i/\n in {0/Z_1, 1/Z_2, 2/Z_3} {
    \pgfmathsetmacro\xx{\i*16}
    \node[zn, minimum width=12mm] (b\i) at (\xx,0) {$\n$};
  }
  \node[zn, tonedk, minimum width=12mm] (zin) at (60,0) {$z$};
  \node[lblit, anchor=west] at (68,0) {incoming};

  % steps
  \node[bxs, anchor=north west, minimum width=124mm, minimum height=6.5mm] (s1) at (-16,-9)
    {\textbf{(1)} packed filter: $\;z \subseteq Z_i$ decided in $O(|G_i|)$ \emph{without} decompressing $Z_i$ \;$\Rightarrow$\; drop};
  \node[bxs, anchor=north west, minimum width=124mm, minimum height=6.5mm] (s2) at (-16,-17)
    {\textbf{(2)} $z \subs Z_i$ for some $i$ \;$\Rightarrow$\; drop $z$};
  \node[bxs, anchor=north west, minimum width=124mm, minimum height=6.5mm] (s3) at (-16,-25)
    {\textbf{(3)} retire every stored $Z_i \subs z$};
  \node[bxs, anchor=north west, minimum width=124mm, minimum height=6.5mm] (s4) at (-16,-33)
    {\textbf{(4)} while some $Z_i$ has $[\,z \sqcup Z_i\,] = [z] \cup [Z_i]$: \; $z := z \sqcup Z_i$, retire $Z_i$};
  \node[bxs, anchor=north west, minimum width=124mm, minimum height=6.5mm] (s5) at (-16,-41)
    {\textbf{(5)} intern $z$, prepend the new node to the chain};

  \node[lbl, anchor=east] at (-6,-52) {\itshape after};
  \node[zn, tonedk, minimum width=12mm] (a0) at (0,-52) {$z'$};
  \node[dead, minimum width=12mm] (a1) at (16,-52) {$Z_1$};
  \node[zn, minimum width=12mm] (a2) at (32,-52) {$Z_2$};
  \node[dead, minimum width=12mm] (a3) at (48,-52) {$Z_3$};
  \node[lblit, anchor=west, text width=44mm] at (60,-52)
    {$z' = z \sqcup Z_1 \sqcup Z_3$;\\dashed $=$ tombstoned};

\end{tikzpicture}
\caption{Insertion into one discrete state's antichain.  Steps (2) and (3)
are the classical subsumption test in both directions.  Step (4) is the
\emph{exact} convex-union merge: if the convex hull of two zones happens to
contain no point outside their union, the two are replaced by the hull with no
loss of precision (\cref{thm:hull}).  Step (4) is what keeps chains at
$\approx 1$ zone per discrete state on every process network we measure, which
in turn is what makes the critical section in \cref{fig:protocol} $O(1)$ rather
than $O(\text{chain})$.  Step (1) is a pure optimisation of step (2) that runs
directly on the compressed representation -- see \cref{fig:mingraph}.}
\label{fig:add}
\end{figure}}
